In the last section, the method by which the \gls{IP}s are found was discussed.
In this section the actual instrumentation is reviewed.
As stated in \cref{sec:instr_point_search}, there are two types of instrumentation:
\begin{itemize}
    \item The first one is used to set up the variables used by the \gls{SHM} update rule.
    \item All the other \gls{IP}s are used for the actual \gls{SHM} update.
\end{itemize}
With reference to \cref{alg:afl_style_update}, there are 3 main variables to consider:
\begin{itemize}
    \item \textit{current\_bb\_id}: ID of the \gls{BB} in which the instrumentation is placed
    \item \textit{previous\_bb\_id}: ID of the \gls{BB} from which the branch is taken
    \item \textit{shared\_map}: Address of the \gls{SHM}
\end{itemize}
In the last example of the previous section (\cref{fig:pow_example}) the \gls{IP}s have already been found.
Without considering the first one (which is where the initial instrumentation is placed), for all the other \gls{IP}s the 
initial address of the \gls{BB} is where the instrumentation will be placed.
Let us consider, for example, the \textbf{0x11f1} \gls{IP}, with previous \gls{BB} equal to \textit{0x11c2} and current \gls{BB} equal to \textit{0x11f1}.
In this case, since there is only one incoming edge to \textbf{0x11f1}, the \textit{previous\_bb\_id} can be placed inline.
But in the general case (e.g. \textbf{0x11b5}) there are multiple edges that end in the \gls{BB}. For this reason, it becomes impossible to 
place the ID of the \gls{BB} directly inline.
The solution implemented in the tool discussed in this thesis consists of storing a variable in a section of the binary.\footnote{
    This section could be the same one used for relocation or a new one added specifically to store the data.
}
This variable represents the \textit{previous\_bb\_id} and is initialized to a random value by the initial instrumentation.

For the \textit{current\_bb\_id}, instead, there is no need to store it in memory. It will be hardcoded (as a random value) in the 
assembly instrumentation inserted within the basic block.

The \textit{shared\_map} has already been resolved and written in the section added to the binary specifically to store 
resolved symbols.

So now the whole instrumentation phase can be reviewed in greater detail.
\begin{center}
    \textbf{Initial Instrumentation}
\end{center}
This instrumentation is needed to:
\begin{itemize}
    \item Initialize the \textit{previous\_bb\_id}. This value is stored in a specific section created for the purpose of 
        holding temporary values needed for instrumentation.
    \item Dereference the resolved \gls{SHM} address. This initialization step is performed for optimization purposes. 
        As can be seen in \cref{fig:shm_relocation}, the instrumented section that works as a \gls{GOT} holds the address 
        of the variable inside the injected \gls{SO}. In order to obtain the actual \gls{SHM} address, the value must be dereferenced twice. 
        If this double dereference had to be performed at every \gls{IP}, it would result in a lot of useless overhead.
        For this reason, in the initial instrumentation the value is dereferenced twice and then saved in the same section as \textit{previous\_bb\_id}.
        By doing this, at each \gls{IP} the \gls{SHM} address only has to be loaded into a register, and the offset of the correct cell computed.
\end{itemize}
\begin{center}
    \textbf{\gls{SHM} update instrumentation}
\end{center}
\cref{alg:shm_update} shows the general approach to the \gls{SHM} update.
Since the value of the \texttt{rax} register is modified, it must be saved before the actual \gls{SHM} update and restored afterwards.
The value of \texttt{rax} is not saved on the stack but in the same section where \textit{previous\_bb\_id} and 
\textit{shared\_map\_address} are stored.
\begin{algorithm}
    \caption{Pseudo-code of Shared Memory update instrumentation}
    \label{alg:shm_update}
    \begin{algorithmic}[1]
        \STATE $\text{Save the state of \texttt{rax} register that will be used for computation}$
        \STATE $\texttt{rax} \gets \textit{previous\_bb\_id}$
        \STATE $\texttt{rax} \gets \texttt{rax} \oplus \textit{current\_bb\_id}$ \COMMENT{\small \textcolor{green!60!black}{\texttt{rax} will contain the result of the hash operation between the basic blocks}}
        \STATE $\texttt{rax} \gets \texttt{rax} + \textit{shared\_map\_address} $\COMMENT{\small \textcolor{green!60!black}{since the shared map is an array, the offset is obtained by adding the initial address to the index}}
        \STATE $[\,\texttt{rax}\,] \gets [\,\texttt{rax}\,] + 1$ \COMMENT{\small \textcolor{green!60!black}{the [\,\,] is the dereference operator}}
        \STATE $\textit{previous\_bb\_id} \gets \textit{current\_bb\_id} \gg 1$
        \STATE $\text{Restore the value of \texttt{rax} register to the previously saved value}$
    \end{algorithmic}
\end{algorithm}

Now the complete approach can be reviewed: 
\begin{itemize}
    \item Add a dynamic dependency on a \gls{SO} that contains the forkserver function.
    \item Add a symbol that will contain the \gls{SHM} address.
    \item Add two new sections. The first will be used as a \gls{GOT} to relocate the symbol. The second will 
        be used as temporary memory space to store \textit{previous\_bb\_id} and \textit{shared\_map\_address}, alongside register saving.
    \item Add a relocation entry in order to resolve the symbol previously added. That symbol is exported from the \gls{SO}.
    \item Find the \gls{IP}s with the algorithm discussed in \cref{sec:instr_point_search}.
    \item Add the initial instrumentation used to initialize the values of 
        \textit{previous\_bb\_id} and \textit{shared\_map\_address} stored in the added section.
    \item For each previously found \gls{IP}, add the \gls{SHM} update instructions as seen in \cref{alg:shm_update}.
\end{itemize}